\subsection{Economic-model boundary and formulation}
After the carbon ledger establishes the denominator for abatement cost, the economic model constructs the numerator. The principal mature case uses the coherent Nishihara reference economics: 59.7 billion JPY plant cost, 0.52 JPY/MJ heat and 4.9 JPY/kWh electricity at 85\% availability~\cite{nishihara2007potential}. Because no validated project off-design electricity output is claimed, the controlling economic test assigns \textbf{S\$0/MWh electricity export value}. To avoid free reactor capacity, the calculation charges the full published source-product burden: the source 370 MWth heat product at its source heat cost plus the source 88 MWe generation at its source electricity cost. This is a conservative economic allocation, not a claim that the project itself exports 88 MWe.

The source's 70.9 billion JPY / 0.57 JPY/MJ / 5.5 JPY/kWh case is retained only as an adverse \emph{cost} sensitivity in which assigned IHX/secondary-loop cost doubles; it is not interpreted as increased exchanger capacity. CCS capital is anchored to IEAGHG Case 1A and T\&S remains a Singapore screening scenario~\cite{ieaghg2017smr,mticcs2025cost}.

Forward abatement cost remains
\begin{equation}\label{eq:abatement-cost-definition}
C_{\rm abatement}=\frac{C_{\rm candidate}-C_{\rm baseline}}{A_{\rm avoided}}.
\end{equation}
where $C_{\rm candidate}$ and $C_{\rm baseline}$ are annual candidate and baseline costs (S\$/y), $A_{\rm avoided}$ is annual lifecycle abatement (tCO$_2$e/y), and $C_{\rm abatement}$ is forward abatement cost (S\$/tCO$_2$e). Section~\ref{sec:ledgers} provides the numerical substitution.

